% ---------------------------------------------------------------------
%  CHƯƠNG 1 — MỞ ĐẦU
% ---------------------------------------------------------------------
\chapter{Mở đầu}
\label{ch:modau}

\begin{muctieu}
Chương mở đầu trình bày lý do chọn đề tài, mục tiêu cần đạt, phạm vi thực
hiện và bố cục báo cáo, giúp người đọc nắm được bức tranh tổng quan về đồ án
xây dựng hệ thống hỗ trợ tìm phòng trọ tại Đà Nẵng.
\end{muctieu}

\section{Lý do chọn đề tài}

Đà Nẵng là một trong những thành phố có số lượng sinh viên và lao động
ngoại tỉnh nhập cư đông, kéo theo nhu cầu thuê phòng trọ rất lớn quanh các
khu vực trường đại học, khu công nghiệp và trung tâm thành phố. Tuy nhiên,
thông tin cho thuê phòng trọ hiện nằm rải rác trên nhiều kênh khác nhau như
mạng xã hội, các nhóm chia sẻ và biển quảng cáo dán trước nhà, thiếu tính
thống nhất và khó tra cứu theo đúng nhu cầu cá nhân. Người tìm trọ thường
phải xem qua hàng chục tin đăng, so sánh thủ công giữa giá thuê, vị trí và
tiện ích, mất nhiều thời gian mà vẫn có thể bỏ sót những lựa chọn phù hợp.
Bài toán gợi ý \en{(recommendation)} trong trí tuệ nhân tạo cho phép tự
động xếp hạng và đề xuất những phòng trọ phù hợp nhất dựa trên tiêu chí của
từng người dùng, thay vì buộc người dùng phải tự lọc qua toàn bộ danh sách.
Đề tài này được chọn vì vừa giải quyết một nhu cầu thực tế gần gũi với đời
sống sinh viên, vừa là cơ hội để em vận dụng kiến thức về hệ gợi ý, xử lý
dữ liệu và xây dựng ứng dụng web đã học trong học phần.

\section{Mục tiêu của đồ án}

Đồ án hướng tới các mục tiêu sau:
\begin{itemize}[leftmargin=1.6em, itemsep=2pt]
  \item Xây dựng được bộ dữ liệu phòng trọ tại khu vực Đà Nẵng với đầy đủ
        thông tin về giá thuê, vị trí, diện tích và tiện ích.
  \item Cài đặt được một mô hình gợi ý phòng trọ dựa trên nội dung, có khả
        năng xếp hạng phòng trọ theo mức độ phù hợp với tiêu chí người
        dùng nhập vào.
  \item Đạt Precision@5 từ 0,7 trở lên khi đánh giá ngoại tuyến trên tập
        tương tác người dùng thử nghiệm.
  \item Phát triển giao diện web cho phép người dùng nhập tiêu chí tìm
        phòng và xem danh sách gợi ý kèm bản đồ vị trí.
  \item So sánh mô hình gợi ý dựa trên nội dung thuần tuý với phiên bản có
        kết hợp tín hiệu từ người dùng tương tự, để làm rõ mức độ cải thiện.
\end{itemize}

\section{Phạm vi và đối tượng}

Đồ án tập trung vào việc gợi ý phòng trọ, nhà trọ cho thuê dài hạn dành cho
sinh viên và người đi làm tại khu vực nội thành Đà Nẵng, cụ thể là sáu
quận trung tâm gồm Hải Châu, Thanh Khê, Sơn Trà, Ngũ Hành Sơn, Cẩm Lệ và
Liên Chiểu. Dữ liệu phòng trọ được thu thập thông qua khảo sát trực tiếp và
từ những tin đăng công khai được chủ trọ đồng ý cung cấp. Đồ án không xử lý
loại hình căn hộ chung cư cao cấp hay khách sạn, không thực hiện chức năng
đặt cọc hay thanh toán trực tuyến, và chưa triển khai ở quy mô sản xuất
phục vụ số lượng người dùng lớn.

\begin{luuy}
Lỗi phổ biến nhất của đồ án là chọn phạm vi quá rộng rồi bỏ dở. Một bài toán
hẹp làm tới nơi tới chốn luôn được chấm cao hơn một ý tưởng lớn còn dang dở.
\end{luuy}

\section{Bố cục báo cáo}

Phần còn lại của báo cáo được tổ chức như sau. Chương~\ref{ch:coso} trình
bày cơ sở lý thuyết và công nghệ sử dụng. Chương~\ref{ch:thietke} phân tích
bài toán và thiết kế giải pháp. Chương~\ref{ch:caidat} mô tả quá trình cài
đặt và thử nghiệm. Chương~\ref{ch:ketqua} trình bày kết quả và đánh giá.
Chương~\ref{ch:ketluan} tổng kết và nêu hướng phát triển.
